\SourceLeaf{350}%
Beziehungen zu unterscheiden sind. Die einen setzen eine bestimmte Gestalt der Fläche im Raum voraus. Die anderen sind unabhängig von den verschiedenen Gestalten, die die Fläche annehmen kann. Um die letzteren geht es hier. Nach dem zuvor Gesagten gehört das Krümmungsmaß dazu. In dieselbe Gruppe fällt, wie man leicht sieht, auch die Untersuchung der auf der Fläche konstruierten Figuren, ihrer Winkel, ihrer Flächeninhalte und ihrer Totalkrümmung, das Verbinden von Punkten durch kürzeste Linien und dergleichen. Alle derartigen Untersuchungen müssen davon ausgehen, dass die Natur der gekrümmten Fläche an sich durch den Ausdruck für ein beliebiges Linienelement in der Form \(\BareRoot(E\mathrm{d}p^2+2F\mathrm{d}p\mathord{.}\mathrm{d}q+G\mathrm{d}q^2)\) gegeben ist. Der Verfasser hat einen Teil seiner seit mehreren Jahren auf diesem Gebiet angestellten Untersuchungen in die vorliegende Abhandlung aufgenommen. Er hat sich auf Untersuchungen beschränkt, die vom Ausgangspunkt nicht zu weit entfernt sind und zum Teil als allgemeine Hilfsmittel für viele weitere Untersuchungen dienen können. In unserer Anzeige müssen wir uns noch stärker beschränken und uns mit einigen Beispielen begnügen. Dazu mögen die folgenden Lehrsätze dienen.

Gehen auf einer gekrümmten Fläche unendlich viele kürzeste Linien gleicher Länge von einem einzigen Anfangspunkt aus, so schneidet die Linie durch ihre Endpunkte jede von ihnen unter einem rechten Winkel. Zieht man in jedem Punkt einer beliebigen Linie auf einer gekrümmten Fläche kürzeste Linien gleicher Länge senkrecht zu dieser Linie, so stehen sie alle auch senkrecht auf der Linie, die ihre anderen Endpunkte verbindet. Beide Lehrsätze, von denen der zweite als Verallgemeinerung des ersten gelten kann, werden sowohl analytisch als auch durch einfache geometrische Überlegungen bewiesen. \textit{Der Überschuss der Winkelsumme eines aus kürzesten Linien gebildeten Dreiecks über zwei rechte Winkel ist gleich der Totalkrümmung des Dreiecks.} Als Winkeleinheit wird dabei der Winkel angenommen, dem ein Bogen von der Länge des Radius entspricht \((57^0\ 17'\ 45'')\). Als Einheit für die ganze Krümmung, aufgefasst als ein Stück der Oberfläche der Hilfskugel, dient der Flächeninhalt eines Quadrats, dessen Seite gleich dem Radius dieser Kugel ist. Dieser wichtige Lehrsatz lässt sich offenbar auch so ausdrücken. Der Überschuss der Winkel eines aus kürzesten Linien gebildeten Dreiecks über zwei rechte Winkel verhält sich zu acht rechten Winkeln wie das Stück der Oberfläche der %
